\documentclass[a4paper]{article}

\usepackage{graphicx}
\usepackage{amsmath,amssymb}
\usepackage{minted}
\usepackage{multirow}
\usepackage{float}
\usepackage{algorithm}
\usepackage{algpseudocode}
\usepackage{todonotes}
\usepackage{forest}
\usepackage{xstring}
\usepackage{xcolor}
\usepackage[most]{tcolorbox}
\usepackage{xparse}
\usepackage{natbib}
\usepackage{adjustbox}
\usepackage[active,tightpage]{preview}

\PreviewBorder=4pt

\usetikzlibrary{graphs,graphdrawing,arrows.meta,positioning,calc,fit,shapes.geometric}
\usegdlibrary{layered}

% for external graphviz
\input{/nd_language_ai_project/src/nd_language_ai/formal/computer/programming/tex/latex/code_clips/external_graphviz.tex}

% for tags
\input{/nd_language_ai_project/src/nd_language_ai/formal/computer/programming/tex/latex/part/control_sequence/kind/semtag/semtag.tex}

% for links
\input{/nd_language_ai_project/src/nd_language_ai/formal/computer/programming/tex/latex/code_clips/seehere.tex}

% for nested lists and sections
\input{/nd_language_ai_project/src/nd_language_ai/formal/computer/programming/tex/latex/code_clips/deep_structure.tex}

\title{}
\date{}
\author{Mohammad Rahmani}

\begin{document}
   \maketitle

   \begin{preview}
      \section*{TikZ live-refresh CPU and floating label editor}

Rapid iPad TikZ changes now update the local vector canvas immediately while canonical source uploads are rate-limited to 20 updates per second. The scheduler always preserves the newest pending source and flushes it before explicit rendering or closing the TikZ workspace.

The responsive bridge no longer launches a TikZ LuaLaTeX renderer after every \texttt{update\_tikz}. Each live update now performs only the atomic source save, desktop update event, and acknowledgement. Initial-open and explicit-render behavior remain available through the existing bridge renderer.

Desktop refresh scheduling now separates three costs. Canonical Source text reloads immediately, Source/Visual graphic projections receive bounded live frames at roughly 80 ms intervals even while a drag continues, and the full PDF preview is compiled only after approximately 650 ms of inactivity. This removes the previous trailing-only debounce that made a continuously moving connector appear frozen until the gesture ended.

Source and Visual now share one TikZ preview coordinator. The conservative named-node and named-connector subset produced by the iPad uses a lightweight Qt vector renderer rather than starting LuaLaTeX, which substantially reduces CPU load during normal direct manipulation. Unsupported complex TikZ falls back to a single serialized LuaLaTeX worker with only the newest pending generation retained; successful older work cannot replace a newer displayed generation.

Node and connector label editing now uses a floating editor at the top of the TikZ canvas instead of an inline field inside the selected object. The selected object's label still previews the draft live, tap-away still commits it, and Enter/Done remains optional. The workspace ignores the keyboard safe-area inset so the iPadOS input assistant may remain visible without pushing the drawing surface upward.

   \end{preview}

   \bibliographystyle{plainnat}
   \bibliography{/home/donkarlo/Dropbox/repo/nd_shared_research/bibliography/bibliography.bib}
\end{document}
